\chapter{Motion Planning}
\label{ch:motion_planning}

Motion planning asks whether an object can move from a start state to a goal
state without colliding with obstacles. The central idea is to replace motion
of a physical object by a path for a point in a higher-dimensional
\emph{configuration space}.

\section{Configuration Space}
\label{sec:configuration_space}

A configuration records all degrees of freedom of a moving object. A point
robot translating in the plane has configuration $q=(x,y)\in\mathbb{R}^2$. A rigid
body translating and rotating in the plane has
\[
q=(x,y,\theta)\in\mathbb{R}^2\times S^1,
\]
while a rigid body in three dimensions has six degrees of freedom,
$q=(x,y,z,\alpha,\beta,\gamma)\in\mathbb{R}^3\times SO(3)$.

The configuration space $\mathcal{C}$ contains all possible configurations.
The obstacle region $\mathcal{C}_{\mathrm{obs}}$ consists of configurations
in which the robot intersects an obstacle. The free configuration space is
\[
\mathcal{C}_{\mathrm{free}} =
\mathcal{C}\setminus\mathcal{C}_{\mathrm{obs}}.
\]
Motion planning then becomes the problem of finding a continuous path
$\tau:[0,1]\to\mathcal{C}_{\mathrm{free}}$ with
$\tau(0)=q_{\mathrm{start}}$ and $\tau(1)=q_{\mathrm{goal}}$.

\section{Configuration-Space Obstacles}
\label{sec:configuration_obstacles}

For a translating robot $R$ and workspace obstacle $O$, collision-free
translations can be computed using the Minkowski difference:
\[
\mathcal{C}_{\mathrm{obs}} = O\oplus(-R).
\]
Thus a finite-sized robot can be replaced by a point, while the obstacles are
expanded by the reflected robot. For a rotating robot, this construction is
performed for each orientation, producing a slice of the higher-dimensional
configuration-space obstacle.

\begin{definition}[Free Path]
A free path is a continuous map $\tau:[0,1]\to\mathcal{C}_{\mathrm{free}}$.
It is \emph{feasible} if it connects the start and goal configurations and
satisfies any required constraints, such as joint limits or velocity bounds.
\end{definition}

\section{Exact Planning Methods}
\label{sec:exact_motion_planning}

\subsection{Cell Decomposition}

Cell decomposition partitions $\mathcal{C}_{\mathrm{free}}$ into simple cells
whose interiors are either entirely free or entirely blocked. A connectivity
graph is built by connecting adjacent free cells. A graph search gives a
sequence of cells, after which a continuous path is constructed through them.
Exact decompositions provide completeness, but their complexity grows rapidly
with the dimension and algebraic complexity of the configuration space.

\subsection{Visibility Graphs}

For a polygonal point robot in the plane, a visibility graph uses the start,
goal, and obstacle vertices as nodes. Two nodes are connected when the segment
between them lies in free space. Dijkstra's algorithm or A* then finds a
shortest collision-free path in the graph. Visibility graphs are exact for
Euclidean shortest paths among polygonal obstacles.

\section{Roadmaps and Sampling-Based Planning}
\label{sec:sampling_motion_planning}

Sampling-based planners avoid explicitly constructing high-dimensional
obstacle boundaries. They sample configurations, reject samples in
$\mathcal{C}_{\mathrm{obs}}$, and connect nearby collision-free samples.

\begin{algorithm}[htbp]
\caption{Probabilistic Roadmap Planner}
\label{alg:probabilistic_roadmap}
\begin{algorithmic}[1]
\Function{BuildRoadmap}{$\mathcal{C}$, $N$}
  \State $G\gets$ empty graph
  \For{$i=1$ to $N$}
    \State Sample $q$ from $\mathcal{C}$
    \If{$q\in\mathcal{C}_{\mathrm{free}}$}
      \State Add $q$ to $G$
      \For{nearby vertex $p$ in $G$}
        \If{the local path from $p$ to $q$ is collision-free}
          \State Add edge $(p,q)$ to $G$
        \EndIf
      \EndFor
    \EndIf
  \EndFor
  \State \Return $G$
\EndFunction
\end{algorithmic}
\end{algorithm}

Probabilistic roadmaps (PRM) are effective for repeated queries in a fixed
environment. Rapidly-exploring random trees (RRT) grow from the start toward
unexplored regions and are useful for single-query planning and systems with
complex dynamics. Both methods are probabilistically complete under standard
sampling and local-planning assumptions, but they do not automatically return
the shortest path.

\section{Planning Constraints and Validation}
\label{sec:motion_planning_validation}

Collision checking must test the entire local path, not only its endpoints.
Practical planners also handle robot joint limits, self-collision, narrow
passages, dynamic obstacles, clearance requirements, and kinodynamic limits on
velocity and acceleration. A returned path should be checked independently,
then shortened or smoothed only with collision tests after every modification.

\section{Applications}

Motion planning is used in robot manipulation, autonomous vehicles, computer
animation, computer games, surgical navigation, warehouse automation, and
molecular modeling. Configuration-space reasoning also connects polygon
offsets, Minkowski operations, shortest paths, visibility, and geometric
collision detection.

